% GENERATED FILE. Edit canonical lessons, then run npm run generate:books.
% canonical-source-hash: fnv1a64:2f071e8fc97a28b2
% canonical-lessons: ML-C473-srroppu, ML-C473-kaikanikkuka, ML-C473-irakkivituka, ML-C473-untuka, ML-C473-otiyettuka

\chapter{Catching the Bus}
\label{ch:ml-catching-the-bus}
\hlchaptermodality{473}

\begin{chapteropening}
\textbf{By the end of this chapter:} \emph{I can say a bus stop, to wave down, to drop off, to shove and to rush up.}

Complete the last lesson of chapter 473: I can say a bus stop, to wave down, to drop off, to shove and to rush up.
\end{chapteropening}

\section[\texorpdfstring{\ml{സ്റ്റോപ്പ്} (sṟṟōppŭ)}{സ്റ്റോപ്പ് (sṟṟōppŭ)}]{\ml{സ്റ്റോപ്പ്} (sṟṟōppŭ) — a bus stop; a stop}

\label{lesson:ML-C473-srroppu}



\begin{quote}
Before the new one: say the Malayalam for to creep along, then the Malayalam for to jolt.
\end{quote}

\subsection*{You'll want to know: \ml{സ്റ്റോപ്പ്}}
\textbf{\ml{സ്റ്റോപ്പ്}} — \emph{sṟṟōppŭ} — \textquotedblleft{}a bus stop; a stop\textquotedblright{}.

\begin{quote}
\ml{അടുത്ത} \ml{സ്റ്റോപ്പിൽ} \ml{ഞാൻ} \ml{ഇറങ്ങും}.
\end{quote}

\begin{quote}
\emph{aṭutta sṟṟōppil ñān iṟaṅṅuṁ.}
\end{quote}

\begin{quote}
I'll get off at the next stop.
\end{quote}

A stop for a bus or a boat; the place where the train halts is a \ml{സ്റ്റേഷൻ}.

\begin{cousinweb}
\textbf{\ml{സ്റ്റോപ്പ്}} (\emph{sṟṟōppŭ}) was borrowed from English \emph{stop}, a place where a bus stops.
\end{cousinweb}

\begin{grammarlens}[title={What you've built}]
One more word for getting around.
\end{grammarlens}

\subsection*{Your turn}
\begin{itemize}
\raggedright
  \item \emph{Say it:} \emph{sṟṟōppŭ}
  \item \emph{Say it:} \emph{sṟṟōppŭ}, once more
  \item \emph{Say it:} \emph{aṭutta sṟṟōppil ñān iṟaṅṅuṁ.} — I'll get off at the next stop.
  \item \emph{Recall:} say the Malayalam for to creep along, then the Malayalam for to jolt, then say \emph{sṟṟōppŭ} again
\end{itemize}

\begin{tcolorbox}[breakable,colback=teal!4,colframe=teal!35!black,title={Before you move on}]
What does \ml{സ്റ്റോപ്പ്} mean? (\textquotedblleft{}A bus stop; a stop\textquotedblright{}.) Say it once more.
\end{tcolorbox}
\section[\texorpdfstring{\ml{കൈകാണിക്കുക} (kaikāṇikkuka)}{കൈകാണിക്കുക (kaikāṇikkuka)}]{\ml{കൈകാണിക്കുക} (kaikāṇikkuka) — to wave down; to signal by hand}

\label{lesson:ML-C473-kaikanikkuka}



\begin{quote}
Before the new one: say the Malayalam for to jolt, then the Malayalam for a bus stop.
\end{quote}

\subsection*{You'll want to know: \ml{കൈകാണിക്കുക}}
\textbf{\ml{കൈകാണിക്കുക}} — \emph{kaikāṇikkuka} — \textquotedblleft{}to wave down; to signal by hand\textquotedblright{}.

\begin{quote}
\ml{ഞാൻ} \ml{കൈകാണിച്ചപ്പോൾ} \ml{ബസ്} \ml{നിർത്തി}.
\end{quote}

\begin{quote}
\emph{ñān kaikāṇiccappōḷ basŭ nirtti.}
\end{quote}

\begin{quote}
When I waved, the bus stopped.
\end{quote}

Past: \ml{കൈകാണിച്ചു}, \emph{kaikāṇiccu}.

Literally to show the hand. On a road in a village you stop a bus this way.

\begin{cousinweb}
\textbf{\ml{കൈകാണിക്കുക}} (\emph{kaikāṇikkuka}) is a native Dravidian compound of \ml{കൈ}, the hand (DEDR 2023), and \ml{കാണിക്കുക}, to show (DEDR 1443).
\end{cousinweb}

\begin{grammarlens}[title={What you've built}]
One more word for getting around.
\end{grammarlens}

\subsection*{Your turn}
\begin{itemize}
\raggedright
  \item \emph{Say it:} \emph{kaikāṇikkuka}
  \item \emph{Say it:} \emph{kaikāṇikkuka}, once more
  \item \emph{Say it:} \emph{ñān kaikāṇiccappōḷ basŭ nirtti.} — When I waved, the bus stopped.
  \item \emph{Recall:} say the Malayalam for to jolt, then the Malayalam for a bus stop, then say \emph{kaikāṇikkuka} again
\end{itemize}

\begin{tcolorbox}[breakable,colback=teal!4,colframe=teal!35!black,title={Before you move on}]
What does \ml{കൈകാണിക്കുക} mean? (\textquotedblleft{}To wave down; to signal by hand\textquotedblright{}.) Say it once more.
\end{tcolorbox}
\section[\texorpdfstring{\ml{ഇറക്കിവിടുക} (iṟakkiviṭuka)}{ഇറക്കിവിടുക (iṟakkiviṭuka)}]{\ml{ഇറക്കിവിടുക} (iṟakkiviṭuka) — to drop off}

\label{lesson:ML-C473-irakkivituka}



\begin{quote}
Before the new one: say the Malayalam for a bus stop, then the Malayalam for to wave down.
\end{quote}

\subsection*{You'll want to know: \ml{ഇറക്കിവിടുക}}
\textbf{\ml{ഇറക്കിവിടുക}} — \emph{iṟakkiviṭuka} — \textquotedblleft{}to drop off\textquotedblright{}.

\begin{quote}
\ml{ഓട്ടോക്കാരൻ} \ml{എന്നെ} \ml{സ്റ്റേഷന്} \ml{മുന്നിൽ} \ml{ഇറക്കിവിട്ടു}.
\end{quote}

\begin{quote}
\emph{ōṭṭōkkāran enne sṟṟēṣanŭ munnil iṟakkiviṭṭu.}
\end{quote}

\begin{quote}
The auto driver dropped me off in front of the station.
\end{quote}

Past: \ml{ഇറക്കിവിട്ടു}, \emph{iṟakkiviṭṭu}.

\ml{ഇറക്കി}, having let down, plus \ml{വിടുക}, to let go. The person dropped takes -\ml{എ}: \ml{എന്നെ}.

\begin{cousinweb}
\textbf{\ml{ഇറക്കിവിടുക}} (\emph{iṟakkiviṭuka}) is a native Dravidian compound of \ml{ഇറക്കുക}, to bring down (DEDR 516), and \ml{വിടുക}, to let go.
\end{cousinweb}

\begin{grammarlens}[title={What you've built}]
One more word for getting around.
\end{grammarlens}

\subsection*{Your turn}
\begin{itemize}
\raggedright
  \item \emph{Say it:} \emph{iṟakkiviṭuka}
  \item \emph{Say it:} \emph{iṟakkiviṭuka}, once more
  \item \emph{Say it:} \emph{ōṭṭōkkāran enne sṟṟēṣanŭ munnil iṟakkiviṭṭu.} — The auto driver dropped me off in front of the station.
  \item \emph{Recall:} say the Malayalam for a bus stop, then the Malayalam for to wave down, then say \emph{iṟakkiviṭuka} again
\end{itemize}

\begin{tcolorbox}[breakable,colback=teal!4,colframe=teal!35!black,title={Before you move on}]
What does \ml{ഇറക്കിവിടുക} mean? (\textquotedblleft{}To drop off\textquotedblright{}.) Say it once more.
\end{tcolorbox}
\section[\texorpdfstring{\ml{ഉന്തുക} (untuka)}{ഉന്തുക (untuka)}]{\ml{ഉന്തുക} (untuka) — to shove; to push}

\label{lesson:ML-C473-untuka}



\begin{quote}
Before the new one: say the Malayalam for to wave down, then the Malayalam for to drop off.
\end{quote}

\subsection*{You'll want to know: \ml{ഉന്തുക}}
\textbf{\ml{ഉന്തുക}} — \emph{untuka} — \textquotedblleft{}to shove; to push\textquotedblright{}.

\begin{quote}
\ml{തിരക്കിൽ} \ml{ആരോ} \ml{എന്നെ} \ml{പിന്നിൽ} \ml{നിന്ന്} \ml{ഉന്തി}.
\end{quote}

\begin{quote}
\emph{tirakkil ārō enne pinnil ninnŭ unti.}
\end{quote}

\begin{quote}
In the crowd someone shoved me from behind.
\end{quote}

Past: \ml{ഉന്തി}, \emph{unti}.

A rougher word than \ml{തള്ളുക}, to push.

\begin{cousinweb}
\textbf{\ml{ഉന്തുക}} (\emph{untuka}) is native Dravidian, \textquotedblleft{}to push, to thrust, to shove\textquotedblright{} (DEDR 625).
\end{cousinweb}

\begin{grammarlens}[title={What you've built}]
One more word for getting around.
\end{grammarlens}

\subsection*{Your turn}
\begin{itemize}
\raggedright
  \item \emph{Say it:} \emph{untuka}
  \item \emph{Say it:} \emph{untuka}, once more
  \item \emph{Say it:} \emph{tirakkil ārō enne pinnil ninnŭ unti.} — In the crowd someone shoved me from behind.
  \item \emph{Recall:} say the Malayalam for to wave down, then the Malayalam for to drop off, then say \emph{untuka} again
\end{itemize}

\begin{tcolorbox}[breakable,colback=teal!4,colframe=teal!35!black,title={Before you move on}]
What does \ml{ഉന്തുക} mean? (\textquotedblleft{}To shove; to push\textquotedblright{}.) Say it once more.
\end{tcolorbox}
\section[\texorpdfstring{\ml{ഓടിയെത്തുക} (ōṭiyettuka)}{ഓടിയെത്തുക (ōṭiyettuka)}]{\ml{ഓടിയെത്തുക} (ōṭiyettuka) — to rush up; to come running}

\label{lesson:ML-C473-otiyettuka}



\begin{quote}
Before the new one: say the Malayalam for to drop off, then the Malayalam for to shove.
\end{quote}

\subsection*{You'll want to know: \ml{ഓടിയെത്തുക}}
\textbf{\ml{ഓടിയെത്തുക}} — \emph{ōṭiyettuka} — \textquotedblleft{}to rush up; to come running\textquotedblright{}.

\begin{quote}
\ml{ബസ്} \ml{പുറപ്പെടുന്നതിന്} \ml{മുൻപ്} \ml{അവൾ} \ml{ഓടിയെത്തി}.
\end{quote}

\begin{quote}
\emph{basŭ puṟappeṭunnatinŭ munpŭ avaḷ ōṭiyetti.}
\end{quote}

\begin{quote}
She came running before the bus left.
\end{quote}

Past: \ml{ഓടിയെത്തി}, \emph{ōṭiyetti}.

\ml{ഓടി}, having run, plus \ml{എത്തുക}, to arrive.

\begin{cousinweb}
\textbf{\ml{ഓടിയെത്തുക}} (\emph{ōṭiyettuka}) is a native Dravidian compound of \ml{ഓടുക}, to run (DEDR 1041), and \ml{എത്തുക}, to reach (DEDR 809).
\end{cousinweb}

\begin{grammarlens}[title={What you've built}]
One more word for getting around.
\end{grammarlens}

\subsection*{Your turn}
\begin{itemize}
\raggedright
  \item \emph{Say it:} \emph{ōṭiyettuka}
  \item \emph{Say it:} \emph{ōṭiyettuka}, once more
  \item \emph{Say it:} \emph{basŭ puṟappeṭunnatinŭ munpŭ avaḷ ōṭiyetti.} — She came running before the bus left.
  \item \emph{Recall:} say the Malayalam for to drop off, then the Malayalam for to shove, then say \emph{ōṭiyettuka} again
\end{itemize}

\begin{tcolorbox}[breakable,colback=teal!4,colframe=teal!35!black,title={Before you move on}]
What does \ml{ഓടിയെത്തുക} mean? (\textquotedblleft{}To rush up; to come running\textquotedblright{}.) Say it once more.
\end{tcolorbox}
